% Separately instantiated corrected French for printed p.122.
% Decision: TOH-P122-R1-SRC-DISPOSITION-001.
% No authorial error is established; historical wording remains unchanged.
\SourcePage{122}

\providecommand{\TohokuContentsChapter}[1]{%
  \par\medskip\noindent\textbf{#1}\par}
\providecommand{\TohokuContentsEntry}[2]{%
  \par\noindent\hangindent=2.8em\hangafter=1%
  \makebox[2.8em][l]{#1}#2\par}

\LayerTarget{tohoku:contents}
\begingroup\small
\LayerTarget{tohoku:section:4.2}
\TohokuContentsEntry{}{fondamentale}
\LayerTarget{tohoku:section:4.3}
\TohokuContentsEntry{4.3.}{Cas d'un faisceau d'anneaux constant}
\LayerTarget{tohoku:section:4.4}
\TohokuContentsEntry{4.4.}{Cas des faisceaux avec groupe d'opérateurs}

\LayerTarget{tohoku:chapter:V}
\TohokuContentsChapter{CHAPITRE V\quad ETUDE COHOMOLOGIQUE DES ESPACES A
OPERATEURS}

\LayerTarget{tohoku:section:5.1}
\TohokuContentsEntry{5.1.}{Généralités sur les $G$-faisceaux}
\LayerTarget{tohoku:section:5.2}
\TohokuContentsEntry{5.2.}{Les foncteurs $H^{n}(X;\,G,A)$ et
$\mathbf{H}^{n}(G,A)$ et les suites spectrales fondamentales}
\LayerTarget{tohoku:section:5.3}
\TohokuContentsEntry{5.3.}{Cas d'un groupe discontinu d'homéomorphismes}
\LayerTarget{tohoku:section:5.4}
\TohokuContentsEntry{5.4.}{Transformation de la première suite spectrale}
\LayerTarget{tohoku:section:5.5}
\TohokuContentsEntry{5.5.}{Calcul des $H^{n}(X;\,G,A)$ par les recouvrements}
\LayerTarget{tohoku:section:5.6}
\TohokuContentsEntry{5.6.}{Les foncteurs
$\mathrm{Ext}_{O,G}^{n}(X;\,A,B)$}
\LayerTarget{tohoku:section:5.7}
\TohokuContentsEntry{5.7.}{Introduction de familles $\Phi$.}
\endgroup

\LayerTarget{tohoku:chapter:I}
\LayerContents{Chapitre I — Généralités sur les catégories abéliennes}
\begin{center}
  \textbf{Chapitre I\quad Généralités sur les catégories abéliennes}
\end{center}

\LayerTarget{tohoku:section:1.1}
\noindent\textbf{1.1. Catégories.}\quad
\LayerTarget{tohoku:section:1.1:p1}
Rappelons qu'on appelle \textit{catégorie} une classe non vide $C$ d'objets,
avec la donnée pour $A,B\in C$ d'un ensemble $\mathrm{Hom}(A,B)$ (appelé
ensemble des \textit{morphismes} de $A$ dans $B$), et pour trois objets
$A,B,C\in C$, d'une application (dite \textit{composition des morphismes})
$(u,v)\longrightarrow vu$ de $\mathrm{Hom}(A,B)\times\mathrm{Hom}(B,C)$ dans
$\mathrm{Hom}(A,C)$, ces données satisfaisant aux deux axiomes suivants : la
composition des homomorphismes est \textit{associative} ; pour tout $A\in C$,
il existe dans $\mathrm{Hom}(A,A)$ un élément $I_A$ (appelé le
\textit{morphisme identique de $A$}) qui soit une unité à droite et à gauche
pour la composition des morphismes. (Cet élément $I_A$ est alors déterminé de
façon unique). Enfin, il sera prudent de supposer que la donnée d'un morphisme
$u$ détermine ses objets de ``départ'' et ``d'arrivée'', en d'autres termes que
si $(A,B)$ et $(A',B')$ sont deux couples distincts d'objets de $C$, alors
$\mathrm{Hom}(A,B)$ et $\mathrm{Hom}(A',B')$ sont deux ensembles disjoints.

\LayerTarget{tohoku:section:1.1:p2}
Si $C$ est une catégorie, on définit la \textit{catégorie duale} $C^{\circ}$
comme la catégorie ayant les mêmes objets que $C$, et où l'ensemble
$\mathrm{Hom}(A,B)^{\circ}$ des morphismes de $A$ dans $B$ est identique à
$\mathrm{Hom}(B,A)$, le composé de $u$ et $v$ dans $C^{\circ}$ étant défini
comme le composé de $v$ et $u$ dans $C$. Toute notion et tout énoncé relatif à
une catégorie non précisée admet une \textit{notion ou un énoncé dual}
(``procédé de renversement des flèches''), qui sera tout aussi utile dans les
applications, mais dont l'explicitation est laissée le plus souvent au
lecteur.

\LayerTarget{tohoku:section:1.1:p3}
Soit donnée une catégorie $C$ et un morphisme $u:A\longrightarrow B$ dans $C$.
Pour tout $C\in C$, on définit une application
$v\longrightarrow uv:\mathrm{Hom}(C,A)\longrightarrow\mathrm{Hom}(C,B)$ et
une application
$w\longrightarrow wu:\mathrm{Hom}(B,C)\longrightarrow\mathrm{Hom}(A,C)$.
On dit que $u$ est un \textit{monomorphisme} ou que $u$ est
\textit{injectif} (resp. que $u$ est un \textit{épimorphisme} ou que $u$ est
\textit{surjectif}) si la première (resp. la seconde) des deux applications
précédentes est toujours injective ; $u$ est dit \textit{bijectif} si $u$ est
à la fois injectif et surjectif.

\LayerTarget{tohoku:section:1.1:p4}
On appelle \textit{inverse à gauche} (resp. \textit{à droite}) de $u$ un
$v\in\mathrm{Hom}(B,A)$ tel que $vu=I_A$ (resp. $uv=I_B$) ; $v$ est dit
\textit{inverse de $u$} s'il est à la fois inverse à gauche et inverse à
droite (auquel cas il est déterminé de façon unique). $u$ est dit un
\textit{isomorphisme} s'il admet un inverse. Si $u$ admet un inverse à gauche
(resp. à droite) il est injectif (resp. surjectif), donc un isomorphisme est

\TohokuContinuation{123}
